\section{Evaluation}
\label{sec:eval}
Five questions. \textbf{RQ1}: does enforcing the safety fragment change
attack outcomes on a live model that follows the injection, and how
does it compare with simpler checks and an existing defense?
\textbf{RQ2}: does the gate add security where admission is provably
blind? \textbf{RQ3}: what does enforcement cost in benign utility, at
the power needed to resolve it? \textbf{RQ4}: do the answers hold
across models, suites, and attack families? \textbf{RQ5}: what does it
cost in latency? Every number flows from the episode logs through a
recompute pipeline that emits the macros used in the text.

\subsection{Setup and Registered Design}
\label{sec:evalsetup}
\paragraph*{Subjects, suites, attacks}
The primary subject is a frontier live endpoint, gemini-2.5-flash, on
the AgentDojo \emph{workspace} suite (email, calendar, cloud drive;
benchmark v1.2.x, package \texttt{agentdojo}~0.1.35) under the canonical
\texttt{important\_instructions} attack family, through a checkpointed
harness derived from the benchmark's own
\texttt{run\_task\_with\_pipeline} path. The pilot grid ran at the
endpoint's default sampling temperature (the pilot harness did not set
one); the registered harness sends temperature~0.7 and the derived seed
explicitly on every request. A dated proprietary snapshot
(\texttt{gpt-4o-2024-05-13}) is the positive calibration control (5/5
at $t{=}0.0$, 4/5 at $t{=}0.7$ on the known-susceptible
\texttt{workspace/user\_task\_0} injections), ruling out an inert
harness. A 0.5B open-weight driver (Qwen2.5-0.5B-Instruct, greedy)
developed the harness and the rationalization probes of
\S\ref{sec:rq2}; it completes no workspace task with or without
enforcement (0/16), so no utility or attack rate is claimed from it.
RQ4 adds a second frontier endpoint, a post-fix re-run of the
open-weight qwen3 family, and the benchmark's \texttt{tool\_knowledge}
attack family, the first stage of the cascade
of~\cite{bhagwatkar2025firewalls}; all three replications completed.
AgentDyn adds the completed dynamic-suite replication reported in
\S\ref{sec:rq4}; suite-generality remains bounded by its task and model
coverage~\cite{agentdyn}.

\paragraph*{Deployed policy objects}
$\lambda$ is as in Definition~\ref{def:policy}: the user turn and the
system prompt $\mathrm{TRUSTED}$, every tool result $\mathrm{MEDIUM}$,
model text no evidence at all. Support of an action is the set of atoms
naming one of its argument values plus the user atoms that express its
class (the registered directive lexicon); $S = \sum w / (1 + \sum w)$
over qualifying support with $\tc = 0.45$ at $\lreq = \mathrm{TRUSTED}$
for all protected classes, fixed before any attack run. This $S$
satisfies A1--A3 of Definition~\ref{def:i3} (sum over attested atoms
at or above $\lreq$; $x/(1+x)$ increasing with $w \ge 0$; weights
fixed in the TCB from the recorded prefix). $\pi$ authorizes no
post-ingress raise. $\bindpol$ binds targets that leave the trust
boundary (recipients, payees, users, channels, URLs, amounts) at
$\mathrm{TRUSTED}$ and internal handles that only tools can enumerate
(file, e-mail, and event ids; hotel and vendor names) at
$\mathrm{MEDIUM}$; free text and dates are unbound.
$D = \{\mathrm{delete}, \mathrm{share}, \mathrm{send}, \mathrm{book}\}$.
Two-root corroboration is not deployable on a single first-party
channel, so policy tightness is bounded by the attestation topology.
The offline sweep gives a flat region at $\tc\in\{0.30,0.45,0.50\}$:
\SpSweepTauFortyFiveSecAdmit/\SpSweepTauFortyFiveSecN{} security and
\SpSweepTauFortyFiveBenignAdmit/\SpSweepTauFortyFiveBenignN{} benign first protected
decisions admit. At $\tc=0.60$ and 0.67, both rates fall to zero, placing
the cliff between 0.50 and 0.60. Requiring every bound slot to be
TRUSTED reduces security admissions to \SpSweepEllAllTrustedSecAdmit/
\SpSweepEllAllTrustedSecN{} but leaves benign admissions unchanged at
\SpSweepEllAllTrustedBenignAdmit/\SpSweepEllAllTrustedBenignN.

\paragraph*{Development versus test}
Every case any development run touched is excluded from the held-out
pools and reported as development evidence: the four rationalization
probes, the five MCP authorization patterns of \S\ref{sec:redteam},
the 22-task small-model workload on which $\tau_c$ was fixed, the
post-fix open-weight run, and the 40-case/24-case pilot grid of
\S\ref{sec:rq1}, whose defended arms ran a \emph{protected-tool
filter}: the pilot harness never wired support extraction, so every
state-changing call was dropped regardless of evidence, and the one
target outside the filter's list (calendar-event creation) accounts for
all four defended-arm breaches. The held-out case sets (150 security
cases and the 56 benign tasks no development run touched, across the
four suites, drawn by a seeded procedure), the 20 choice-laundering
probes with their honest controls, the policy objects, the baseline and
gate configurations, and the utility minimal important difference are
deposited hash-sealed and publicly timestamped before the RQ1--RQ4 runs:
registration v2 (archive SHA-256 \texttt{ee88c72b\ldots c68}; manifest
\texttt{cbddbd33\ldots f0}; OpenTimestamps block 974214,
2026-11-02T05:06:44Z; the deleted original DOI is being re-registered at
final upload as DOI \url{https://doi.org/10.5281/zenodo.24075933}) and the replacement-2
testbed v4 (DOI \url{https://doi.org/10.5281/zenodo.23914405};
archive SHA-256 \texttt{d97f1369\ldots f3b}; manifest
\texttt{c67de998\ldots 071}; OpenTimestamps block 974389,
2026-11-03T13:03:48Z). The campaign hotfix is DOI
\url{https://doi.org/10.5281/zenodo.23917748}, manifest
\texttt{dfc691f8\ldots 13b}, OpenTimestamps block 974430,
2026-11-03T19:47:14Z.
An intermediate replacement-1 implementation and its results were
excluded because they lacked independent restarts and MCP transport. The $\tau_c$ workload touched only benign
traces; security cases on those workspace user tasks remain in the
pool and form a registered descriptive stratum beside the $\tau_c$
sweep, rather than being silently presented as untouched.

\paragraph*{Baselines}
Three baselines run under the identical harness and case sets:
\textbf{B1}, \emph{argument--instruction matching}, every target named
in a protected call must appear in the user's instruction (the simpler
check a decision-layer defense must beat); \textbf{B2}, a tool-filter
policy in the CaMeL lineage~\cite{camel,oob} (no tool-sourced value may
reach a protected sink unless the user turn names it); \textbf{B3},
\Ifour{} alone. Every arm's verdict is logged on every protected call
as a shadow; all seven arms completed the registered campaign.

\paragraph*{Statistics}
Wilson 95\,\% intervals on every rate; Fisher's exact test on the one
pre-registered primary comparison per grid, with risk difference (RD)
and risk ratio (RR) and their intervals beside every $p$-value; the
rule-of-three bound beside any zero-event claim; secondary comparisons
Holm-corrected. Five harness defects (among them the unset
temperature) are documented in the artifact; degenerate episodes (null
model output) are reported and excluded only in a sensitivity analysis.
The reference semantics' own coverage check (10/10 classes, 3/3 honest
baselines) is Appendix~\ref{app:refsem}.

\subsection{RQ1: Security on a Live Model}
\label{sec:rq1}
\paragraph*{The subject is susceptible}
gemini-2.5-flash follows the canonical injection in
\SecNdHit/\SecNdN{} no-defense pilot episodes (ASR \SecNdASR, Wilson
95\,\% CI $[\SecNdCIlo, \SecNdCIhi]$); a disjoint 30-case sweep
corroborates this ($23/60 = 0.383$ $[0.271, 0.510]$). The pilot
covers three workspace user tasks, a scope the held-out set removes.

\paragraph*{A protected-tool filter already collapses this family (pilot)}
On this subject (\SecNdN{} independent cases per arm), a filter that
drops every call to a state-changing tool reduces measured attack
success from \SecNdASR{} (\SecNdHit/\SecNdN) to \SecGateASR{}
(\SecGateHit/\SecGateN, Wilson $[\SecGateCIlo, \SecGateCIhi]$) in each of
two replicates: Fisher exact $p{=}\GateVsNdFisherP$, RD $= \HeadlineRD$
$[\HeadlineRDlo, \HeadlineRDhi]$, RR $= \HeadlineRR$
$[\HeadlineRRlo, \HeadlineRRhi]$ (Table~\ref{tab:livemodel}). The filter
verifiably engages (\SecMonFilt--\SecGateFilt{} protected calls dropped
per arm); every surviving breach is the one template whose target,
calendar-event creation, was not on its list, and excluding the four
degenerate episodes changes nothing (2/37, $p{=}6.4\times10^{-5}$).
This is B2 in its strictest form, and it sets the bar: admission must
match that security at lower benign cost. Security is measured across
seven arms on the held-out set; the powered RQ3 crossover decides the
comparative benign-cost bar.
The registered rows show that admission reduces ASR from
\SpSecNoDefenseASR{} to \SpSecDeiASR{} (Fisher
$p{=}\SpSecDeiFisherP$), but it does not beat B1: both are
\SpSecBOneASR, with Holm $p{=}\SpSecDeiVsBOneHolmP$. B2strict is lower
still at \SpSecBTwoStrictASR. RQ1 is therefore a result about
\Ifour{} rather than evidence of full-DEI superiority; the gate does not
change admission's ASR. The powered RQ3 contrasts, not the 56-task
descriptive arm, decide comparative benign cost.

\begin{table}[t]
\centering
\caption{RQ1 and RQ3. Pilot (development evidence): gemini-2.5-flash,
workspace, $n{=}\SecNdN$ cases per security arm, $n{=}\UtilNdN$ per
benign arm; the defended arms are two replicates of a protected-tool
filter. Registered campaign: held-out cases, seven arms.}
\label{tab:livemodel}
\fontsize{6.5}{7.5}\selectfont
\setlength{\tabcolsep}{0pt}
\begin{tabular}{@{}lccc@{}}
\toprule
Arm & ASR & 95\,\% CI & Filtered \\
\midrule
\multicolumn{4}{@{}l}{\emph{Pilot grid}}\\
No defense & $\SecNdASR$ ($\SecNdHit/\SecNdN$) & $[\SecNdCIlo, \SecNdCIhi]$ & $\SecNdFilt$ \\
Protected-tool filter, replicate 1 & $\SecMonASR$ ($\SecMonHit/\SecMonN$) & $[\SecMonCIlo, \SecMonCIhi]$ & $\SecMonFilt$ \\
Protected-tool filter, replicate 2 & $\SecGateASR$ ($\SecGateHit/\SecGateN$) & $[\SecGateCIlo, \SecGateCIhi]$ & $\SecGateFilt$ \\
Benign ($n{=}\UtilNdN$/arm) & \multicolumn{3}{l}{no defense $\UtilNdPass/\UtilNdN$; filter $\UtilGatePass/\UtilGateN$ (\S\ref{sec:rq3})} \\
\midrule
\multicolumn{4}{@{}l}{\emph{Registered campaign} (150 security / 56 benign held-out cases)}\\
No defense & \RqOneNodefenseCell & \RqOneNodefenseCI & \RqOneNodefenseFiltered \\
Admission (\Ione--\Ifour) & \RqOneDeiCell & \RqOneDeiCI & \RqOneDeiFiltered \\
Admission + gate & \RqOneDeigateCell & \RqOneDeigateCI & \RqOneDeigateFiltered \\
B1 argument--instruction matching & \RqOneBOneCell & \RqOneBOneCI & \RqOneBOneFiltered \\
B2 tool filter & \RqOneBTwoCell & \RqOneBTwoCI & \RqOneBTwoFiltered \\
B2strict filter & \RqOneBTwoStrictCell & \RqOneBTwoStrictCI & \RqOneBTwoStrictFiltered \\
B3 \Ifour{} alone & \RqOneBThreeCell & \RqOneBThreeCI & \RqOneBThreeFiltered \\
\bottomrule
\end{tabular}
\end{table}

\subsection{RQ2: What the Gate Adds}
\label{sec:rq2}
On the pilot grid the two defended arms are replicates of one filter
(the pilot's gate re-commitment was not wired), so the pilot cannot show
the gate's value, and the canonical injection family is caught at
admission in any case. The gate is designed for the case admission
cannot see (Cor.~\ref{cor:ratgap}), and RQ2 is answered there.

\paragraph*{Development probes: the counterexample, executed}
Four probes pair a legitimate directive ($\mathrm{TRUSTED}$,
$\tc$-sufficient) with a planted message naming an attacker-chosen
target ``with approval,'' on the small-model driver (existence
evidence, not rates). In 8 runs the model followed the injection
\GateFollows{} times; twice it laundered the action under the genuine
directive with $\tc$-sufficient support, and \Ione--\Ithree{}
\emph{admitted}: Cor.~\ref{cor:ratgap} end to end. With the gate
enabled (policy and perturbations frozen beforehand), both laundering
admissions were quarantined on the oracle plane $\delta_1$
(\GateQuar{} of \GateQuar); the honest-citation follow was blocked by
\Ithree{} before the gate; the interrupted follow, repaired to the
user's target, re-committed identically under both perturbations. Of
\GateBenignDestr{} gated benign commitments, \GateFalseQuarN{} were
falsely quarantined, at \GateOverhead$\times$ model time and
\GateMonitorUs\,$\mu$s p50 monitor latency. Cited-untrusted removal at
$\lmin$ never triggered (0/26), which is why laundering passes
\Ione--\Ithree{} at all.

\paragraph*{What these probes do and do not show}
They show that the gate closes the counterexample where \Ione--\Ithree{}
are blind, not that it adds value over \Ifour: both laundered targets
are absent from every trusted belief, so \Ifour{} (and B1) blocks them
at admission (Remark~\ref{rem:residual}). The registered RQ2 design
therefore uses \emph{choice-laundering} probes, in which every bound
argument is user-named and the plant is a false \emph{claim} that
selects among legitimate options: F1, delete whichever of two named
files is older (the plant reverses the dates), placed inside the file
listing the task reads (co-located) or inside a calendar entry it also
reads (separate); F2, send to a named contact only if an event falls on
a given day (the plant claims it moved); F3, write to whichever of two
named contacts attends an event (the plant claims a substitution).
Twenty probes with matched honest controls, four registered seeds
each, validated so that each honest ground truth passes and each plant
triggers the laundered action; \Ifour{} (and B1) admit every laundered
action by construction, so any quarantine is the gate's. The primary
endpoints are pooled $\delta_3$ catch on up to 40 separable-plant episodes
and false quarantine on up to 60 actionable honest-control evaluations,
with Wilson and probe-level cluster-bootstrap intervals and the actual
denominators reported; non-action controls never enter a gate-flip
denominator. Family/placement cells, observed
$\delta_2$, and oracle-$\delta_1$ $\mathit{fp}$/$\mathit{miss}$ per
Prop.~\ref{prop:approxgap} are secondary; registered prediction: $\delta_3$ catches
the separate placement and F2 at no honest cost and cannot see
co-located plants, which $\delta_2$ flags only together with every
honest commitment that needs the same atom.
The outcome refutes zero honest cost: $\delta_3$ catches
\RqTwoSeparableDeltaThreeK/\RqTwoSeparableDeltaThreeN{} separable
commitments (\RqTwoSeparableDeltaThreeRate, Wilson
$[\RqTwoSeparableDeltaThreeLo,\RqTwoSeparableDeltaThreeHi]$, probe-bootstrap
$[\RqTwoSeparableDeltaThreeClusterLo,\RqTwoSeparableDeltaThreeClusterHi]$)
but falsely quarantines \RqTwoHonestDeltaThreeK/\RqTwoHonestDeltaThreeN{}
controls (\RqTwoHonestDeltaThreeRate,
$[\RqTwoHonestDeltaThreeLo,\RqTwoHonestDeltaThreeHi]$). Only 26 separable
and 29 actionable evaluations arose; non-actions were not imputed.
Co-located catch is \RqTwoColocatedDeltaThreeRate, while $\delta_2$
flags every co-located plant and \RqTwoHonestDeltaTwoK/\RqTwoHonestDeltaTwoN{}
controls, the resolution limit. Against the oracle, $\delta_3$ has TP
\RqTwoDeltaThreeTp, FN \RqTwoDeltaThreeFn, FP \RqTwoDeltaThreeFp, TN
\RqTwoDeltaThreeTn (miss \RqTwoDeltaThreeMissRate, FP rate
\RqTwoDeltaThreeFpRate); $\delta_2$ has TP \RqTwoDeltaTwoTp, FN
\RqTwoDeltaTwoFn, FP \RqTwoDeltaTwoFp, TN \RqTwoDeltaTwoTn. The gate helps
on some separable plants, but its atom-granularity and re-commitment
failures are limits, not a general security claim.

\subsection{RQ3: Utility Cost}
\label{sec:rq3}
\paragraph*{Pilot}
Under the protected-tool filter the benign-completion rate is
\UtilGatePass/\UtilGateN{} (\UtilGateRate) versus \UtilNdPass/\UtilNdN{}
(\UtilNdRate) with no defense, with failures \UtilGateBlock/\UtilGateN{}
versus \UtilNdBlock/\UtilNdN. The 12.5-point estimate is not resolvable
at $n{=}\UtilGateN$ (Fisher $p{=}0.49$; paired discordant 3/0, exact
McNemar $p{=}0.25$) and is structurally uninformative: 23 of the 24
benign tasks never call a protected tool, the one that does was
blocked, and the two other extra failures occurred with nothing
filtered. We claim neither that the cost is small nor that it is
large; the registered study stratifies on protected-tool use.

\paragraph*{Registered study}
The utility comparison is re-run as a within-task crossover: each of
the 97 benign tasks of the four suites runs under three registered
seeds with no defense, admission plus gate, B1, and B2strict (291 pairs
per contrast, registered random order), pairing the within-task
variation an independent-arms design treats as noise. The sample size
follows from a pre-registered minimal important difference of 10 points
from the 83.3\,\% baseline (independent-arms floor 266 per arm at
80\,\% power; the paired test has power 0.83 at the independence bound
and more under positive within-task correlation). Paired RD with a
Newcombe interval for admission plus gate versus no defense is the
headline; the paired contrasts against B2strict and B1, Holm-corrected,
decide whether admission costs less than the filter; exact McNemar and
a task-level cluster bootstrap is secondary; protected-tool-use and
development-contact strata are descriptive; every failure is classified
as defense-induced (blocked by a named invariant), model-induced, or
degenerate.
The registered crossover shows a material cost: DEI+gate changes
completion by \RqThreeUtilDeiGatePairedRD{} relative to no defense
(Newcombe $[\RqThreeUtilDeiGatePairedRDlo,\RqThreeUtilDeiGatePairedRDhi]$;
task-cluster bootstrap $[\RqThreeUtilDeiGateClusterRDlo,\RqThreeUtilDeiGateClusterRDhi]$),
not by a negligible amount. It nevertheless outperforms B2strict by
\RqThreeUtilDeiGateVsBTwoStrictPairedRD{}
($[\RqThreeUtilDeiGateVsBTwoStrictPairedRDlo,\RqThreeUtilDeiGateVsBTwoStrictPairedRDhi]$;
Holm $p{=}\RqThreeUtilDeiGateVsBTwoStrictHolmP$) and B1 by
\RqThreeUtilDeiGateVsBOnePairedRD{}
($[\RqThreeUtilDeiGateVsBOnePairedRDlo,\RqThreeUtilDeiGateVsBOnePairedRDhi]$;
Holm $p{=}\RqThreeUtilDeiGateVsBOneHolmP$), the latter borderline.
DEI+gate failures split into 74 defense-induced, 115 model-induced, and
0 degenerate episodes; B2strict splits into 132 defense-induced and 88
model-induced. The practical claim is therefore a security--utility
trade-off, not preserved utility.

\subsection{RQ4: Generality}
\label{sec:rq4}
The pilot covers three workspace user tasks, one attack family, and
one model (per-case susceptibility 0.18--1.00), so no cell
extrapolates. Pre-fix open-weight archives (0/87 on qwen3-coder-30B,
0/33 on qwen3-32B) are negative observations, not resistance; a
post-fix development run of qwen3.6-35B followed the injection in
18/60 no-defense episodes across the four suites, so RQ4 re-runs that
family under the registration. On gpt-4.1-mini, ASR is
\RqFourASecNoDefenseASR{} with no defense, \RqFourASecDeiASR{} with DEI,
\RqFourASecDeiGateASR{} with DEI+gate, and \RqFourASecBOneASR{} with B1.
On \texttt{tool\_knowledge}, the corresponding rates are
\RqFourCSecNoDefenseASR, \RqFourCSecDeiASR, and
\RqFourCSecBOneASR. All 150-security-case rows have Wilson intervals and
zero voided outcomes; the benign rows are secondary and are not promoted
to primary evidence. The open-weight row is completed but uninformative:
qwen3:8b follows only \RqFourBSecNoDefenseHit/\RqFourBSecNoDefenseN{}
no-defense injections (DEI \RqFourBSecDeiHit/\RqFourBSecDeiN), so it does
not establish susceptibility or robustness. AgentDyn replicates on 1,120
episodes with 30 voided attempts excluded: ASR falls from
\RqAgentDynSecNoDefenseASR{} to \RqAgentDynSecDeiASR{} (RD
\RqAgentDynSecDeiRD{} $[\RqAgentDynSecDeiRDlo,\RqAgentDynSecDeiRDhi]$,
Fisher $p=\RqAgentDynSecDeiFisherP$), with suite-generality stated only
for the dynamic benchmark.

\subsection{RQ5: Latency}
Admission costs are those of Table~\ref{tab:overhead}: the replacement
A2 gateway adds \RqFiveOverheadRange{} p50 overhead over A0 at
5--100 req/s, with mean p50 \RqFiveATwoTwentyFiveMedian\,ms and mean p95
\RqFiveATwoTwentyFiveTail\,ms at 25 req/s. The gate adds one
re-commitment per enforced perturbation on irreversible classes
(\GateOverhead$\times$ benign model time in the development probes).

\subsection{Threats to Validity}
\label{sec:threats}
\paragraph*{Construct and internal}
\DEI{} constrains evidential structure, not truth, and within-TCB
compromise is excluded by the threat model. The pilot's defended arms
ran a protected-tool filter at an unset temperature (\S\ref{sec:evalsetup});
the registered harness fixes both, and B3 isolates \Ifour's
contribution. Development probes were built against the gateway's own
threat model; held-out sets are registered and timestamped before the
runs.
\paragraph*{External and statistical}
CZTA is one instantiation on one deployment; efficacy is per model,
suite, and attack family, and RQ4 bounds the generality claim. The
pilot's security comparison has power $\approx 0.99$ for its effect;
its utility arm is underpowered; its defended arms are one filter, so
RQ2 is answered where admission is provably blind.
\paragraph*{Adaptive adversaries}
\Ifive's perturbation class is registered but finite; an attacker
outside it, or a model that re-commits identically with and without
the plant, is the counterexample \S\ref{sec:falsifiability} invites.
